% from ega4-11-fr.tex
\oldpage[IV-3]{168}
des noyaux des
$(u_\lambda)_\mathfrak p:M_\mathfrak p\to(N_\lambda)_\mathfrak p$ déduits
des $(u_\lambda)$ par le changement de base
$\operatorname{Spec}(B_\mathfrak p)\to\operatorname{Spec}(B)$ est réduite à
0. Autrement dit, on est ramené au cas où $B$ est un anneau local noethérien,
et en considérant l'idéal premier de $A$ image réciproque de l'idéal maximal de
$B$, on peut aussi supposer que $A$ est un anneau local, d'idéal maximal
$\mathfrak m$. Alors, comme $\mathfrak mB$ est contenu dans l'idéal maximal de
$B$, et que $M$ est un $B$-module de type fini, l'intersection des
$\mathfrak m^nM$ est réduite à 0 $(\mathbf 0_{\mathrm I},7.3.5)$, donc il
suffit de prouver que pour tout $n$, l'image de $t$ dans
$M/\mathfrak m^{n+1}M$ est nulle. Il suffit donc de prouver que la famille
déduite de $(u_\lambda)$ par le changement de base
$\operatorname{Spec}(B/\mathfrak m^{n+1}B)\to\operatorname{Spec}(B)$ est
séparante, ce qui signifie encore qu'on peut se borner au cas où $A$ est un
anneau local dont l'idéal maximal $\mathfrak m$ est nilpotent. Mais alors les
$\mathfrak m^k/\mathfrak m^{k+1}$ sont des $(A/\mathfrak m)$-modules libres,
et en vertu de l'hypothèse sur les $(u_\lambda)_s$, on est précisément dans les
conditions d'application de (11.9.15), d'où la conclusion annoncée.

Soient maintenant $h:S'\to S$ un morphisme changement de base, et
$(u'_\lambda)$ la famille obtenue à partir de $(u_\lambda)$ par le changement
de base $h':X\times_S S'\to X$~; prouvons que $(u'_\lambda)$ est aussi
séparante. Supposons d'abord que $h$ soit \emph{localement de type fini}~; il
en est alors de même de $h'$, donc $X'=X\times_S S'$ est localement
noethérien~; de plus, si $s'\in S'$ est au-dessus du point $s\in S$, il résulte
de (11.9.13) appliqué à $X_s$ et à
$X_{s'}=X_s\otimes_{k(s)}k(s')$ que, pour tout $s'\in S'$, la famille
$((u'_\lambda)_{s'})$ est séparante~; on peut par suite conclure de la première
partie de la démonstration que dans ce cas $(u'_\lambda)$ est séparante.

Enfin, si $h$ est quelconque, on peut évidemment se limiter au cas où
$S=\operatorname{Spec}(A)$ et $S'=\operatorname{Spec}(A')$ sont affines, et
considérer $A'$ comme limite inductive de ses sous-$A$-algèbres de type fini.
Comme les $\mathcal G_\lambda$ sont $S$-plats, il suffit d'appliquer ce qui
précède et le lemme (11.9.5) pour terminer la démonstration.
\end{proof}

\begin{proposition}[11.9.17]
\label{IV.11.9.17-fr}
Soient $f:X\to S$ un morphisme localement de présentation finie, $\mathcal F$
un $\mathcal O_X$-Module quasi-cohérent de présentation finie et $f$-plat,
$U$ un ouvert de $X$, $j:U\to X$ l'injection canonique,
$u:\mathcal F\to j_*(j^*(\mathcal F))$ l'homomorphisme canonique
$(\mathbf 0_{\mathrm I},4.4.3.2)$. Pour tout $s\in S$, soient $X_s$ la fibre
$X\times_S\operatorname{Spec}(k(s))$, $U_s$ l'ouvert $U\cap X_s$ de $X_s$,
$\mathcal F_s=\mathcal F\otimes_{\mathcal O_S}k(s)$,
$j_s:U_s\to X_s$ l'injection canonique,
$u_s:\mathcal F_s\to(j_s)_*((j_s)^*(\mathcal F_s))$ l'homomorphisme
canonique. Pour que $u$ soit universellement injectif relativement à $S$, il
faut et il suffit que $u_s$ soit injectif pour tout $s\in S$.
\end{proposition}

\begin{proof}
Il n'y a encore à prouver que la suffisance de la condition. Lorsque $X$ est
localement noethérien, la proposition est un corollaire immédiat de (11.9.16).
Nous allons nous ramener à ce cas en deux étapes, en nous bornant, comme on
peut évidemment le faire, au cas où $S=\operatorname{Spec}(A)$ et $X$ sont
affines.

\emph{A) Cas où $U$ est quasi-compact.} --- Nous utiliserons le lemme suivant~:

\begin{lemma}[11.9.17.1]
\label{IV.11.9.17.1-fr}
Sous les hypothèses générales de (11.9.17), et en supposant en outre $S$ et
$X$ affines et $U$ quasi-compact, l'ensemble $E$ des $s\in S$ tels que $u_s$
soit injectif est constructible.
\end{lemma}

\begin{proof}
En effet, les fibres $X_s$ sont des préschémas localement noethériens, donc
$E$ peut aussi, en vertu de (5.10.2), être défini comme l'ensemble des
$s\in S$ tels que $\operatorname{Ass}(\mathcal F_s)\subset U_s$. Notons en
outre que $U$, étant quasi-compact dans un schéma affine, est constructible.
Alors, la vérification de la condition (9.2.1, (i)) découle aussitôt de
(4.2.7)~; d'autre
